\section{State Trie Structure}
\label{sec:trie}

\subsection{Merkle Patricia Trie}

Lux C-Chain uses a modified Merkle Patricia Trie (MPT) identical to Ethereum's \cite{wood2014ethereum}:

\begin{definition}[Merkle Patricia Trie]
An MPT is a deterministic key-value store where:
\begin{itemize}
    \item Keys are 256-bit hashes (account addresses hashed via Keccak-256).
    \item Values are RLP-encoded account states or storage slots.
    \item Internal nodes are branch nodes (16 children + value), extension nodes (shared prefix + child), or leaf nodes (key remainder + value).
    \item Each node's identity is its Keccak-256 hash (32 bytes).
    \item The trie root hash commits to the entire state.
\end{itemize}
\end{definition}

\begin{definition}[State Root]
The state root $R = H(\text{root node})$ is a 32-byte hash that uniquely determines the entire state trie. Block headers include $R$, making it part of consensus.
\end{definition}

\subsection{Trie Statistics for Lux C-Chain}

\begin{table}[H]
\centering
\caption{Lux C-Chain state trie statistics (2025).}
\label{tab:stats}
\begin{tabular}{lr}
\toprule
\textbf{Metric} & \textbf{Value} \\
\midrule
Total accounts & 48,000,000 \\
Total storage slots & 320,000,000 \\
Trie nodes (account trie) & 95,000,000 \\
Trie nodes (storage tries) & 640,000,000 \\
Total state size (raw) & 42 GB \\
Total state size (trie nodes) & 78 GB \\
Trie depth (average) & 7.2 \\
Trie depth (maximum) & 12 \\
\bottomrule
\end{tabular}
\end{table}

